\begin{tikzpicture}[x=1cm,y=-.87cm,every node/.style={font=\fontsize{12.2}{15}\selectfont,text=Ink},arr/.style={-{Stealth[length=2mm]},line width=.8pt,draw=Classical}]\path[use as bounding box] (0,0) rectangle (7.45,6.2);
\draw[draw=Classical,fill=Classical!8,rounded corners=2pt,line width=.7pt] (0.4,1.3) rectangle (2.9,3.4);
\node[anchor=center,text=Ink,font=\fontsize{16}{15}\selectfont] at (1.65,2.35) {$\mathrm dy/\mathrm dz$};
\draw[arr,draw=Quantum] (3.1,2.35)--(4.1,2.35);
\node[anchor=center,text=Ink,font=\fontsize{11}{15}\selectfont] at (3.6,1.4) {× r(z)};
\draw[draw=Quantum,fill=Quantum!8,rounded corners=2pt,line width=.7pt] (4.3,1.3) rectangle (6.8,3.4);
\node[anchor=center,text=Ink,font=\fontsize{16}{15}\selectfont] at (5.55,2.35) {$\mathrm dy/\mathrm dt$};
\node[anchor=center,text=Ink,font=\fontsize{12.2}{15}\selectfont,text width=6.7cm,align=center] at (3.7,4.7) {Il tasso locale collega le due unità.};
\node[anchor=center,text=Ink,font=\fontsize{10.8}{15}\selectfont,text width=6.8cm,align=center] at (3.7,5.7) {La relazione vale nei tratti regolari della cronologia.};

\end{tikzpicture}